\documentclass[
  spanish,
  letterpaper, oneside
]{article}

% PRODUCTO: Reporte de investigación (técnica 76 del catálogo UnADM).
% Ficha oficial: López, A. (2022), 100 Técnicas Didácticas, Fascículo 4, pp. 173-183.
% Los cinco elementos de la ficha se alojan en el esqueleto de tres actos: (1) y (5)
% coinciden con Introducción y Conclusión; (2), (3) y (4) son subsecciones del desarrollo.
% Planeación S7, Tema 3: Bloque de regularidad constitucional (CT 293/2011).
% Fuentes: 9 tesis SCJN en referencias-.../notas-.../actividad-6-reporte-de-investigacion.

% Los metadatos de portada NO admiten \pendiente: init.tex los evalúa con
% \ifthenelse antes de que la macro exista.
% El subtítulo viaja al pie de página: máximo 68 caracteres o choca con el curso.
\def\documenttitle {Diferencias entre el control de constitucionalidad y el de convencionalidad}
\def\documentsubtitle {Reporte de lectura sobre criterios jurisdiccionales de la SCJN}
\def\documentsubject {Actividad 7 - Garantías Constitucionales}

\def\documentauthor {Martin Jonathan de la Cruz}
\def\coursename {Garantías Constitucionales}
\def\coursecode {025}

\def\universityname {Universidad Abierta y a Distancia de México}
\def\universityfaculty {Licenciatura en Derecho}
\def\universitydepartment {Garantías Constitucionales}
\def\universitydepartmentimage {departamentos/UnADM}
\def\universitydepartmentimagecfg {height=1.57cm}
\def\universitylocation {Roma Norte, Ciudad de México}

\def\authortable {
  \begin{tabular}{ll}
    \textbf{Alumno:}
    & \begin{tabular}[t]{l}
      Martin Jonathan de la Cruz \\
    \end{tabular} \\
    Matrícula: & ES2611202040 \\
    Figura docente: & Aarón López Pérez \\
    Semestre/Bloque: & 2 / 1 \\
    & \\
    \multicolumn{2}{l}{Fecha de realización: \today} \\
    \multicolumn{2}{l}{\universitylocation}
  \end{tabular}
}

\input{template}
\usepackage{helvet}
\renewcommand{\familydefault}{\sfdefault}
\usepackage{array}
\usepackage{booktabs}
\usepackage{longtable}
\usepackage{pdflscape}
\usepackage{tikz}
\usetikzlibrary{arrows.meta,positioning,calc,fit,shapes.geometric}
\setcitestyle{authoryear,open={(},close={)}}

% Plantilla-Informe no define \pendiente (sí lo hacen las plantillas maestras).
% Marca en rojo lo que falta redactar; al completar la actividad no debe quedar ninguno.
\providecommand{\pendiente}[1]{\textcolor{red}{[PENDIENTE: #1]}}

\def\coverwatermarkenabled {true}
\def\coverwatermarkimage {img/departamentos/UnADM.pdf}
\def\coverwatermarkopacity {0.16}
\def\coverwatermarkwidth {0.72\paperwidth}
\def\coverwatermarkxshift {0cm}
\def\coverwatermarkyshift {-0.35cm}

\newcommand{\insertcoverwatermark}{%
  \ifthenelse{\equal{\coverwatermarkenabled}{true}}{%
    \AddToShipoutPictureBG*{%
      \begin{tikzpicture}[remember picture,overlay]
        \node[
          opacity=\coverwatermarkopacity,
          inner sep=0pt,
          xshift=\coverwatermarkxshift,
          yshift=\coverwatermarkyshift
        ] at (current page.center) {%
          \includegraphics[width=\coverwatermarkwidth]{\coverwatermarkimage}%
        };
      \end{tikzpicture}%
    }%
  }{}%
}

\begin{document}

\insertcoverwatermark
\templatePortrait
\templatePagecfg
\onehalfspacing

% Pasos 1, 2 y 9 de la ficha (portada, índice y referencias) los resuelve la
% plantilla: no se rehacen dentro del cuerpo.
\templateIndex
\templateFinalcfg

%% ===========================================================================
%% ACTO 1 - Introducción.
%% Cubre el elemento (1) de la ficha: problema, antecedentes, hipótesis y
%% objetivos. Es la MISMA sección de introducción del reporte, no una aparte.
%% ===========================================================================
\section{Introducción}

La reforma constitucional de junio de 2011 dejó al juez mexicano con dos herramientas
que se parecen y no son lo mismo. Una le permite contrastar cualquier norma con la
Constitución; la otra, con los tratados internacionales de derechos humanos y con la
jurisprudencia que los interpreta. Ambas pueden ejercerse de oficio, ambas terminan
en la inaplicación de la norma y ambas se invocan con frecuencia en la misma
sentencia. Esa proximidad explica que se confundan, y la confusión no es inocua: de
ella dependen la competencia del órgano que resuelve, el efecto de su decisión y el
parámetro que debe emplear para justificarla.

El presente reporte de lectura examina esa distinción a partir de los criterios
jurisdiccionales asignados y de la bibliografía obligada de la unidad. La pregunta
que lo ordena es la siguiente: \textbf{si el control de constitucionalidad y el de
convencionalidad comparten la técnica de la inaplicación, ¿en qué consiste realmente
su diferencia y qué consecuencias prácticas se siguen de ella?} Sostengo que la
respuesta no está en el procedimiento, sino en el parámetro de contraste y en el
alcance del efecto; y que, tras la Contradicción de Tesis 293/2011, la diferencia
tiende a disolverse en un único control de regularidad, aunque subsiste una
asimetría relevante que conviene no perder de vista.

\section{Dos controles frente a una misma norma}

\subsection{El parámetro de contraste como criterio de distinción}

La lectura de los criterios asignados permite fijar primero aquello que separa a
ambos controles. El \textbf{control de constitucionalidad} confronta la norma o el
acto con la Constitución y con la interpretación que de ella hace la Suprema Corte.
El \textbf{control de convencionalidad} la confronta con los tratados
internacionales de derechos humanos y con la jurisprudencia de la Corte
Interamericana de Derechos Humanos \citep{scjnTesis2001605}. La diferencia inicial,
por tanto, es de \textbf{parámetro}: cambia la norma que sirve de medida, no la
operación que se realiza sobre la norma medida.

Esta precisión importa porque el vocabulario de la reforma puede inducir a error. Al
reconocer el artículo 1.\textsuperscript{o} los derechos humanos de fuente
constitucional y convencional, podría suponerse que todo juez quedó habilitado para
resolver cualquier conflicto normativo invocando indistintamente uno u otro cuerpo.
\textbf{No obstante}, la Corte precisó que la obligación se ejerce siempre
\textit{«en el ámbito de sus competencias»} \citep{scjnTesis2005056}, lo que
significa que el parámetro se amplió, pero la competencia no.

\subsection{Concentrado y difuso: la vía y el efecto}

La segunda distinción no corre entre constitucionalidad y convencionalidad, sino
dentro de cada una: la que media entre control concentrado y control difuso. El
\textbf{control concentrado} corresponde a los órganos del Poder Judicial de la
Federación mediante vías específicas —amparo, controversia constitucional y acción
de inconstitucionalidad— y puede culminar en una declaratoria de invalidez con
efectos generales. El \textbf{control difuso} habilita a cualquier autoridad
jurisdiccional para dejar de aplicar una norma en el caso concreto, sin expulsarla
del ordenamiento \citep{scjnTesis2005720}.

La consecuencia práctica es considerable y la Corte la enunció sin ambigüedad: el
ejercicio del control difuso \textit{no} implica declaratoria general de
inconstitucionalidad \citep{scjnTesis2005057}. El juez que inaplica resuelve su
asunto; la norma sigue vigente para los demás. \textbf{Por ello}, quien confunde
ambos planos suele pedir al juzgador ordinario un pronunciamiento que éste no puede
emitir, o atribuir a una inaplicación un alcance derogatorio que no posee.

Conviene añadir que los dos controles no se excluyen ni se ordenan jerárquicamente
entre sí. El ejercicio del control difuso \textit{«no limita ni condiciona»} el del
concentrado \citep{scjnTesis2002264}, y tampoco opera como remedio subsidiario al
que sólo se acuda cuando falta otra vía \citep{scjnTesis2010143}. Son cauces
paralelos con efectos distintos, no etapas sucesivas de un mismo procedimiento.

\subsection{La obligación ex officio y sus condiciones}

Un tercer rasgo comparte ambos controles y merece atención propia: su carácter
oficioso. Los órganos jurisdiccionales deben ejercerlos aunque las partes no lo
soliciten \citep{scjnTesis2003679}, criterio que se extendió también a la
jurisdicción contencioso-administrativa \citep{scjnTesis2005116}. La razón es
comprensible: si la protección de los derechos humanos dependiera de que el
justiciable supiera invocarla con la técnica adecuada, quedaría condicionada a la
calidad de su defensa.

Ahora bien, «de oficio» no equivale a «sin método». La Corte fijó condiciones
generales para su ejercicio \citep{scjnTesis2005622} y estableció una secuencia
ordenada que debe agotarse en este orden \citep{scjnTesis2000073}:

\begin{enumerate}
  \item \textbf{Interpretación conforme en sentido amplio:} leer la norma a la luz de
        la Constitución y de los tratados, favoreciendo en todo tiempo la protección
        más amplia de la persona.
  \item \textbf{Interpretación conforme en sentido estricto:} cuando caben varias
        lecturas válidas, preferir aquella que resulte acorde con el parámetro de
        regularidad constitucional.
  \item \textbf{Inaplicación:} sólo procede cuando las alternativas anteriores se
        agotaron sin lograr la compatibilidad, y su efecto se limita al caso concreto.
\end{enumerate}

Además, cada órgano lo ejerce dentro de su competencia y
no más allá de ella \citep{scjnTesis2009179}. \textbf{Es decir}, la oficiosidad
describe el deber de examinar, no una licencia para inaplicar sin justificación.

\subsection{El límite: restricciones constitucionales expresas}

Aquí aparece la asimetría que impide tratar a ambos controles como equivalentes
perfectos. La jurisprudencia de la Corte Interamericana resulta vinculante para los
jueces mexicanos \citep{scjnTesis2005942}; sin embargo, el control de
convencionalidad no autoriza a desconocer las restricciones expresas que la propia
Constitución establece \citep{scjnTesis2006165}. Es la misma solución que fijó la
Contradicción de Tesis 293/2011 al construir el parámetro de regularidad
constitucional \citep{scjnTesis2006224, scjnIijCT293y2017}.

\textbf{A mi juicio}, este es el punto donde la distinción conserva verdadera
utilidad práctica. Si los derechos de ambas fuentes integran un mismo parámetro, la
pregunta por cuál control se ejerce pierde relevancia en la mayoría de los casos:
el juez contrasta con un bloque unificado. \textbf{Sin embargo}, cuando la
Constitución contiene una restricción expresa, el origen de la norma vuelve a
importar, porque la disposición constitucional prevalece. La diferencia, entonces,
no desapareció: quedó reservada para los casos límite, que son precisamente los más
difíciles.

\section{Conclusión}

El examen de los criterios asignados permite sostener tres conclusiones. La primera
es que la diferencia entre control de constitucionalidad y de convencionalidad no
radica en la técnica —ambos pueden ejercerse de oficio y ambos culminan en la
inaplicación— sino en el \textbf{parámetro de contraste}: la Constitución en un
caso, los tratados de derechos humanos y la jurisprudencia interamericana en el
otro.

La segunda es que la distinción operativamente decisiva no corre entre ambos
controles, sino entre sus modalidades concentrada y difusa. De ahí dependen la
competencia del órgano y el efecto de la resolución: invalidez general en un caso,
inaplicación al caso concreto en el otro. \textbf{Considero} que buena parte de los
errores de planteamiento en la práctica forense provienen de confundir estos dos
ejes, que son independientes entre sí.

La tercera es que, tras la Contradicción de Tesis 293/2011, ambos controles
convergen en un control único de regularidad constitucional, con la salvedad de las
restricciones constitucionales expresas. \textbf{Desde mi perspectiva}, esa
convergencia es coherente con el artículo 1.\textsuperscript{o} y favorece la
protección efectiva de la persona, porque libera al juzgador de una discusión sobre
etiquetas y lo obliga a razonar sobre el contenido del derecho invocado. La reserva
de las restricciones expresas, en cambio, mantiene abierta una tensión con el
estándar interamericano que la jurisprudencia todavía administra caso por caso, y
que difícilmente podrá resolverse sin una decisión de fondo sobre la relación entre
soberanía constitucional y compromiso convencional.\footnote{En la elaboración de
este trabajo se utilizaron herramientas de inteligencia artificial como apoyo para
la organización y revisión del texto, conforme a los Lineamientos para el uso ético
y responsable de la Inteligencia Artificial en el ámbito académico de la UnADM. El
análisis, la selección de fuentes y la postura sostenida son de mi autoría.}

\bibliography{garantias-constitucionales}

\end{document}
